% TODO make hw stylesheet
\documentclass[11pt]{scrartcl}
\usepackage[a4paper, total={6in, 8in}]{geometry}
\usepackage[usenames,dvipsnames,svgnames]{xcolor}
\usepackage[shortlabels]{enumitem}
\usepackage[framemethod=TikZ]{mdframed}
\usepackage{amsmath,amssymb,amsthm}
\usepackage[colorlinks]{hyperref}
\usepackage{multicol}
\usepackage{tabularx}

\hypersetup{
	colorlinks=true,
	linkcolor=blue,
	filecolor=magenta,      
	urlcolor=magenta,
	pdftitle={MAT 320 HW}
}

\usepackage{microtype}
\usepackage{mathtools}
\usepackage[headsepline]{scrlayer-scrpage}
\usepackage{thmtools}
\usepackage[textsize=footnotesize]{todonotes}

\mdfdefinestyle{mdbluebox}{roundcorner=10pt,innerbottommargin=9pt,
	linecolor=blue,backgroundcolor=TealBlue!5,}
\declaretheoremstyle[headfont=\sffamily\bfseries\color{MidnightBlue},
mdframed={style=mdbluebox},]{thmbluebox}

\mdfdefinestyle{mdbloobox}{frametitlefont=\bfseries,innerbottommargin=8pt,
	nobreak=true,backgroundcolor=TealBlue!5,linecolor=blue,}
\declaretheoremstyle[headfont=\bfseries\color{MidnightBlue},
mdframed={style=mdbloobox},headpunct={\\[3pt]},postheadspace=0pt,]{thmbloobox}

\mdfdefinestyle{mdredbox}{frametitlefont=\bfseries,innerbottommargin=8pt,
	nobreak=true,backgroundcolor=Salmon!5,linecolor=RawSienna,}
\declaretheoremstyle[headfont=\bfseries\color{RawSienna},
mdframed={style=mdredbox},headpunct={\\[3pt]},postheadspace=0pt,]{thmredbox}

\mdfdefinestyle{mdgreenbox}{linecolor=ForestGreen,backgroundcolor=ForestGreen!5,
	linewidth=2pt,rightline=false,leftline=true,topline=false,bottomline=false,}
\declaretheoremstyle[headfont=\bfseries\sffamily\color{ForestGreen!70!black},
mdframed={style=mdgreenbox},headpunct={ --- },]{thmgreenbox}

\mdfdefinestyle{mdorangebox}{linecolor=RedOrange,backgroundcolor=RedOrange!5,
	linewidth=2pt,rightline=false,leftline=true,topline=false,bottomline=false,}
\declaretheoremstyle[headfont=\bfseries\sffamily\color{RedOrange!70!black},
mdframed={style=mdorangebox},headpunct={ --- },]{thmorangebox}

\mdfdefinestyle{mdblackbox}{linecolor=black,backgroundcolor=RedViolet!5!gray!5,
	linewidth=3pt,rightline=false,leftline=true,topline=false,bottomline=false,}
\declaretheoremstyle[mdframed={style=mdblackbox}]{thmblackbox}

\declaretheoremstyle[spaceabove=6pt,spacebelow=6pt]{basehead}

\declaretheorem[style=basehead,name=Problem]{problem}
\declaretheorem[style=thmredbox,name=Problem,numbered=no]{problem*}
\declaretheorem[style=thmbloobox,name=Setup]{setup} % for config geo
\declaretheorem[style=thmredbox,name=Example,sibling=problem]{example}
\declaretheorem[style=thmredbox,name=$\bigstar$ Problem,sibling=problem]{niceprob}
\declaretheorem[style=thmbluebox,name=Theorem,sibling=problem]{theorem}
\declaretheorem[style=thmbluebox,name=Corollary,sibling=theorem]{corollary}
\declaretheorem[style=thmbluebox,name=Proposition,sibling=theorem]{prop}
\declaretheorem[style=thmbluebox,name=Lemma,numberwithin=problem]{lemma}
\declaretheorem[style=thmbluebox,name=Theorem,numbered=no]{theorem*}
\declaretheorem[style=thmbluebox,name=Lemma,numbered=no]{lemma*}
\declaretheorem[style=thmgreenbox,name=Claim,numberwithin=problem]{claim}
\declaretheorem[style=thmgreenbox,name=Claim,numbered=no]{claim*}
\declaretheorem[style=thmblackbox,name=Remark,numberwithin=problem]{remark}
\declaretheorem[style=thmblackbox,name=Remark,numbered=no]{remark*}
\declaretheorem[style=thmgreenbox,name=Definition,sibling=theorem]{definition}
\declaretheorem[style=thmgreenbox,name=Definition,numbered=no]{definition*}
\declaretheorem[style=thmorangebox,name=Warning,sibling=theorem]{warning}
\declaretheorem[style=thmorangebox,name=Warning,numbered=no]{warning*}
\declaretheorem[style=thmorangebox,name=Note,numbered=no]{note}
\declaretheorem[style=thmblackbox,name=Exercise,sibling=problem]{exercise}
\declaretheorem[style=thmblackbox,name=Exercise,numbered=no]{exercise*}
\declaretheorem[style=thmblackbox,name=Question,sibling=problem]{question}
\declaretheorem[style=thmblackbox,name=Question,numbered=no]{question*}

\addtolength{\textheight}{3.14cm}
\providecommand{\ol}{\overline}
\providecommand{\ul}{\underline}
\providecommand{\eps}{\varepsilon}
\providecommand{\half}{\frac{1}{2}}
\providecommand{\dang}{\measuredangle} %% Directed angle
\providecommand{\CC}{\mathbb C}
\providecommand{\FF}{\mathbb F}
\providecommand{\NN}{\mathbb N}
\providecommand{\QQ}{\mathbb Q}
\providecommand{\RR}{\mathbb R}
\providecommand{\ZZ}{\mathbb Z}
\providecommand{\dg}{^\circ}
\providecommand{\ii}{\item}
\providecommand{\setdiff}{\smallsetminus}
\providecommand{\alert}[1]{{\sffamily\textbf{\textcolor{blue}{#1}}}}
\providecommand{\inv}{^{-1}}
\providecommand{\mb}[1]{\mathbf{#1}}

\reversemarginpar

\newenvironment{soln}{\begin{proof}[Solution]}{\end{proof}}
\newenvironment{walkthrough}{\noindent \textbf{Walkthrough.}}{\medskip}
\newlist{walk}{enumerate}{3}
\setlist[walk]{label=\bfseries (\alph*)}

\providecommand{\pow}{\operatorname{Pow}}
\providecommand{\vocab}[1]{{\textbf{\textcolor{ForestGreen}{#1}}}}
\providecommand{\lcm}{\operatorname{lcm}}
\providecommand{\abs}[1]{\left|#1\right|}

\usepackage{asymptote}
\begin{asydef}
	defaultpen(fontsize(10pt));
	unitsize(1cm);
	size(8cm); // set a reasonable default
	usepackage("amsmath");
	usepackage("amssymb");
	settings.tex="pdflatex";
	settings.outformat="pdf";
	// Replacement for olympiad+cse5 which is not standard
	import geometry;
	// recalibrate fill and filldraw for conics
	void filldraw(picture pic = currentpicture, conic g, pen fillpen=defaultpen, pen drawpen=defaultpen)
	{ filldraw(pic, (path) g, fillpen, drawpen); }
	void fill(picture pic = currentpicture, conic g, pen p=defaultpen)
	{ filldraw(pic, (path) g, p); }
	// some geometry
	pair foot(pair P, pair A, pair B) { return foot(triangle(A,B,P).VC); }
	pair orthocenter(pair A, pair B, pair C) { return orthocentercenter(A,B,C); }
	pair centroid(pair A, pair B, pair C) { return (A+B+C)/3; }
	// cse5 abbreviations
	path CP(pair P, pair A) { return circle(P, abs(A-P)); }
	path CR(pair P, real r) { return circle(P, r); }
	pair IP(path p, path q) { return intersectionpoints(p,q)[0]; }
	pair OP(path p, path q) { return intersectionpoints(p,q)[1]; }
	path Line(pair A, pair B, real a=0.6, real b=a) { return (a*(A-B)+A)--(b*(B-A)+B); }
	// cse5 more useful functions
	picture CC() {
		picture p=rotate(0)*currentpicture;
		currentpicture.erase();
		return p;
	}
	pair MP(Label s, pair A, pair B = plain.S, pen p = defaultpen) {
		Label L = s;
		L.s = "$"+s.s+"$";
		label(L, A, B, p);
		return A;
	}
	pair Drawing(Label s = "", pair A, pair B = plain.S, pen p = defaultpen) {
		dot(MP(s, A, B, p), p);
		return A;
	}
	path Drawing(path g, pen p = defaultpen, arrowbar ar = None) {
		draw(g, p, ar);
		return g;
	}
\end{asydef}

\usepackage{mathrsfs}

\ihead{\footnotesize MAT 320 HW10}
\ohead{\footnotesize November 27, 2024}

\title{\vspace{-2em}MAT 320 HW10}
\author{\large Aditya Pahuja}
\date{\large Due 11/27}
\begin{document}
\maketitle

\begin{problem}
	Explain with the $\eps$-$\delta$ definition why
	$f(x) = |x|$ is continuous at $x = 0$ but not differentiable.
\end{problem}

\begin{soln}
	Let $\eps > 0$; then, setting $\delta := \eps$ gives
	\[|x - 0| = |x| < \delta = \eps \implies ||x| - |0|| = ||x|| = |x| < \eps,\]
	which is what it means to be continuous at $0$.
	On the other hand, for $f$ to be differentiable at zero,
	we need to check the existence of the limit
	\[\lim_{x\to0} \frac{f(x) - f(0)}{x - 0}
	= \lim_{x\to0} \frac{|x|}{x}.\]
	In particular, to show that the limit doesn't exist,
	it suffices to find two sequences $(a_n)$ and $(b_n)$
	that both converge to zero but for which
	\[\lim_{n\to\infty} \frac{f(a_n) - f(0)}{a_n - 0}
	\ne \lim_{n\to\infty} \frac{f(b_n) - f(0)}{b_n - 0}.\]
	We do this by taking $a_n = \frac 1n$ and $b_n = -\frac 1n$, since
	the two limits then evaluate to
	\[\lim_{n\to\infty} \frac{|a_n|}{a_n} =
	\lim_{n\to\infty} \frac{|1/n|}{1/n} = \lim_{n\to\infty}1 = 1,\quad
	\lim_{n\to\infty} \frac{|b_n|}{b_n} =
	\lim_{n\to\infty} \frac{|-1/n|}{-1/n} = \lim_{n\to\infty}-1 = -1,\]
	and of course $-1\ne 1$. Thus the limit that defines the derivative
	doesn't exist at zero, so $f$ is not differentiable at zero.
\end{soln}

\begin{problem}
	Show that the function $f\colon\RR\to\RR$ defined by
	\[f(x) = \begin{cases}
		0 & \text{if }x = 0\\
		x^2\sin(\frac1x)&\text{otherwise}
	\end{cases}\]
	is differentiable but not of class $C^1$.
\end{problem}

\begin{soln}
	Clearly the function is differentiable when $x\ne 0$
	because $x^2$, $\sin x$, and $\frac 1x$ are
	all differentiable functions of $x$, thus so are their compositions
	and products, whenever the resulting expression is defined.
	To check that $f$ is differentiable at zero, we need
	to show the existence of
	\[\lim_{x\to 0} \frac{f(x) - f(0)}{x - 0}
	= \lim_{x\to 0} \frac{x^2\sin(\frac1x)}{x}.\]
	This limit is equal to
	\[\lim_{x\to0} x\cdot\sin\left(\frac 1x\right),\]
	which is well-known to be $1$, so $f'(0)$ exists and is equal to $1$.
	
	On the other hand, the derivative when $x\ne 0$ evaluates to
	\[2x\cdot \sin\left(\frac 1x\right)
	- x^2\cdot\cos\left(\frac1x\right)\cdot\left(-\frac{1}{x^2}\right)
	= 2x\sin\left(\frac1x\right) - \cos\left(\frac1x\right).\]
	We wish to show that the derivative is not continuous on $\RR$,
	so we will show that the limit of $2x\sin(\frac1x) - \cos(\frac1x)$
	as $x$ approaches zero is not $f'(0) = 1$.
	Indeed, the sequence $a_n = \frac{1}{2\pi n}$ converges to zero, and
	\[\lim_{n\to\infty} 2a_n \sin\left(\frac 1{a_n}\right)
	- \cos\left(\frac1{a_n}\right)
	= \lim_{n\to\infty} \frac{2}{2\pi n}\sin(2\pi n) - \cos(2\pi n)
	= \lim_{n\to\infty} 0 - 1 = -1.\]
	Thus, if the limit of $f'(x)$ exists as $x\to 0$,
	it must equal $-1\ne f'(0)$, so $f'$ is not a continuous function on $\RR$.
	(In fact the limit doesn't exist at all because the $\cos(1/x)$
	term oscillates way too much.)
\end{soln}

\begin{problem}
	Consider the function $f\colon\RR\to\RR$ defined by
	\[f(x) = \begin{cases}
		x^2 & \text{if }x\ge 0\\
		0 &\text{otherwise}.
	\end{cases}\]
	Perform the following tasks:
	\begin{enumerate}[(a)]
		\ii Show that $f$ is differentiable at $x = 0$.
		\ii Calculate the derivative function $f'\colon\RR\to\RR$.
		\ii Is $f'$ continuous? Is it differentiable?
	\end{enumerate}
\end{problem}

\begin{soln}
	For (a), we need to show the existence of the limit
	\[\lim_{x\to0} \frac{f(x) - f(0)}{x - 0}.\]
	For any $\eps > 0$, take $\delta := \eps$. Then,
	if $|x - 0| < \delta$, we consider two cases:
	\begin{itemize}
		\ii If $-\delta < 0 < x$, then $\frac{f(x) - f(0)}{x - 0} = 0 < \eps$.
		\ii If $0\le x < \delta$, then
		$\frac{f(x) - f(0)}{x - 0} = \frac{x^2}{x} = x < \delta = \eps$.
	\end{itemize}
	In either case,
	\[|x - 0| < \delta \implies \abs{\frac{f(x) - f(0)}{x - 0} - 0} < \eps,\]
	which is exactly what it means for
	\[\lim_{x\to 0} \frac{f(x) - f(0)}{x - 0} = 0,\]
	hence the limit exists (and is equal to zero).
	
	We can also compute $f'(x) = 0$ when $x < 0$
	because the restriction of $f$ to $(-\infty, 0)$ is just the zero function
	and $f'(x) = 2x$ when $x > 0$
	because the restriction of $f$ to the open interval $(0,\infty)$
	is $x\mapsto x^2$. Thus, to answer (b),
	\[f'(x) = \begin{cases}
		2x & x > 0\\
		0 & x \le 0.
	\end{cases}\]
	
	Finally, we see that $f'(x)$ can be written instead as
	\[f'(x) = |x| + x\]
	(since when $x > 0$, $|x| + x = x + x = 2x$, while
	when $x \le 0$, $|x| + x = -x + x = 0$).
	This is a sum of continuous functions, hence continuous,
	but it is not differentiable, since $f'$ being differentiable
	would mean $f'(x) - x = |x|$ is a sum of differentiable functions,
	hence differentiable --- but we showed in Problem 1 that
	$|x|$ is not differentiable at zero.
\end{soln}

\setcounter{problem}{4}

\begin{problem}
	Evaluate the series
	\[\sum_{n=1}^\infty \frac{n}{2^n}\]
	and
	\[\sum_{n=1}^\infty \frac{3^n}{n!}.\]
\end{problem}

\begin{soln}
	For the first series, we start with the geometric series
	\[\sum_{n=0}^\infty x^n,\]
	which converges to $\frac{1}{1-x}$ when $|x| < 1$.
	By the theory we developed in class, the derivative of this function is
	\[\sum_{n=1}^\infty nx^{n-1} = \frac{1}{(1-x)^2}\]
	on the same interval $|x| < 1$.
	We can then multiply this identity by $x$ to get the identity
	\[\sum_{n=1}^\infty nx^n = \frac{x}{(1-x)^2}\]
	whenever $|x| < 1$. Since $-1 < \half < 1$, we can therefore compute
	\[\sum_{n=1}^\infty n\cdot\frac{1}{2^n} =
	\frac{1/2}{(1 - 1/2)^2} = \frac{1/2}{1/4} = \boxed 2.\]
	
	For the second series, we notice that
	\[1 + \sum_{n=1}^\infty \frac{3^n}{n!} =
	\sum_{n=0}^\infty \frac{3^n}{n!} = e^3,\]
	so the series evaluates to $\boxed{e^3 - 1}$.
\end{soln}

\setcounter{problem}{6}

\begin{problem}
	Let $f$ and $g$ be two differentiable functions $\RR\to\RR$.
	Show that, given $a < b$, there is a $c\in(a, b)$ such that
	\[[f(b) - f(a)]\cdot g'(c) = [g(b) - g(a)]\cdot f'(c).\]
	Explain how this statement, Rolle's theorem, and the mean value theorem
	relate to each other.
\end{problem}

\begin{soln}
	Consider the function $h\colon\RR\to\RR$ given by
	\[h(x) := (f(b) - f(a))g(x) - (g(b) - g(a))f(x).\]
	This is a linear combination of differentiable functions,
	so it is differentiable. Further, we can compute
	\[h(a) = h(b) = f(b)g(a) - f(a)g(b),\]
	so Rolle's theorem tells us that there is a point in $c\in (a, b)$
	at which $h'(c) = 0$, which is exactly what we want.
	
	This theorem is a generalization of the mean value theorem
	(which is in turn a generalization of Rolle's theorem),
	since setting $g(x) := x$ for all $x$ spits out the mean value theorem,
	as there exists a $c\in(a, b)$ such that
	\[f(b) - f(a) = (f(b) - f(a))\cdot g'(c) = (g(b) - g(a))\cdot f'(c)
	= (b - a)\cdot f'(c).\qedhere\]
\end{soln}

\begin{problem}
	Let $f\colon[0,1]\to\RR$ be a differentiable function
	such that $f(0) = 0$ and $f(x) > 0$ for all $x\in(0,1)$.
	Prove that there is a number $c\in(0,1)$ such that
	\[\frac{2f'(c)}{f(c)} = \frac{f'(1-c)}{f(1-c)}.\]
\end{problem}

\begin{soln}
	Consider the function $g(x) = f(x)^2\cdot f(1-x)$.
	Its derivative is equal to
	\[2f(x)\cdot f'(x)\cdot f(1-x) - f(x)^2\cdot f'(1-x).\]
	Moreover, $g(0) = g(1) = 0$, so there is a $c\in(0,1)$
	for which $g'(c) = 0$. Since $f(c) > 0$, we can conclude that
	\[0 = \frac{g(c)}{f(c)}
	2f'(c)f(1-c) - f(c)f'(1-c),\]
	which rearranges to
	\[2f'(c)f(1-c) = f(c)f'(1-c) \iff
	\frac{2f'(c)}{f(c)} = \frac{f'(1-c)}{f(1-c)}.\qedhere\]
\end{soln}


























\end{document}